% Organización de la materia
\section{Organización de la materia}\label{cap:01}\label{sec:organizacion}

El programa oficial recorre una cadena conceptual que comienza con transformaciones lineales
y termina en problemas geodésicos y cartográficos:
\[
\boxed{\text{álgebra lineal}}
\longrightarrow
\boxed{\text{geometría diferencial}}
\longrightarrow
\boxed{\text{superficies y elipsoide}}
\longrightarrow
\boxed{\text{proyecciones y marcos}}
\]

\begin{center}
\begin{tabularx}{\textwidth}{>{\bfseries}p{34mm} X p{30mm}}
\toprule
Bloque & Contenido & Estado en este apunte\\
\midrule
Transformaciones &
Rotaciones, reflexiones, traslaciones, homotecias, similaridades, afinidades y
transformaciones proyectivas &
Desarrollado\\
Herramientas lineales &
Matrices, composición, producto interno y vectorial, determinante, producto exterior,
núcleo, formas lineales, puntos fijos, autovalores y autovectores, exponencial &
Desarrollado\\
Cambio de base y SVD &
Bases de autovectores, potencias de una matriz, valores singulares, inversa y
pseudoinversa &
Desarrollado (cap.~\ref{cap:svd}; clase 7 dictada el 7/10, sin transcripción)\\
Coordenadas y variedades &
Coordenadas curvilíneas, bases locales, espacios tangentes y derivadas &
Desarrollado (cap.~\ref{cap:derivada}; clases 10--14 anticipadas respecto de 2C2026)\\
Superficies &
Métrica, esfera, elipsoide, curvatura y geodésicas &
Desarrollado (cap.~\ref{cap:elipsoide}; anticipado, pendiente de validar)\\
Aplicaciones &
Proyecciones conformes, marcos de referencia, Molodensky, grillas y remuestreo &
Panorama conceptual\\
Transformaciones conformes y proyecciones (cierre) &
Cuatro caracterizaciones de la conformidad, Tissot, latitud isométrica, Mercator,
autálica y planteo de Gauss--Krüger &
Desarrollado (cap.~\ref{cap:conformes}; complementario)\\
\bottomrule
\end{tabularx}
\end{center}

\begin{atencion}[Estado del documento]
Los capítulos iniciales siguen el material efectivamente recibido hasta ahora. Los temas
posteriores se presentan como mapa del programa y deberán ampliarse cuando la cátedra entregue
las clases y guías correspondientes. De este modo no se confunde una anticipación conceptual
con contenido ya dictado.
\end{atencion}

\section*{Programa}
\addcontentsline{toc}{section}{Programa}

Cada clase del dictado tiene su capítulo y su unidad interactiva; los complementos (C1 a C4)
tratan temas sin clase propia.

{\small
\begin{longtable}{@{}>{\raggedright\arraybackslash}p{9mm}>{\raggedright\arraybackslash}p{34mm}>{\raggedright\arraybackslash}p{44mm}>{\raggedright\arraybackslash}p{42mm}>{\raggedright\arraybackslash}p{22mm}@{}}
\toprule
\textbf{Clase} & \textbf{Tema} & \textbf{Capítulo} & \textbf{Unidad interactiva} & \textbf{Observaciones}\\
\midrule
\endhead
\bottomrule
\endlastfoot
1 & Rotaciones & \ref{cap:rotaciones}. \nameref{cap:rotaciones} & \linkUnidad{00_desaprender_el_cbc}{Fundamentos}, \linkUnidad{01a_punto_quieto}{Rotaciones I}, \linkUnidad{01b_rotacion_plano}{Rotaciones II}, \linkUnidad{01c_rotaciones_3d}{Rotaciones III}, \linkUnidad{02_producto_interno}{Producto interno} & \\
2 & Más rotaciones, composición, líneas coordenadas & \ref{cap:traslaciones}. \nameref{cap:traslaciones} & \linkUnidad{01c_rotaciones_3d}{Rotaciones III}, \linkUnidad{03_reflexiones_similaridades}{Reflexiones y similaridades} & \\
3 & Transformaciones afines & \ref{cap:afines}. \nameref{cap:afines} & \linkUnidad{03_reflexiones_similaridades}{Reflexiones y similaridades}, \linkUnidad{04_afines_homogeneas}{Afinidades y homogéneas} & \\
4 & Espacios lineales y transformaciones lineales & \ref{cap:formas-lineales}. \nameref{cap:formas-lineales} & \linkUnidad{00_desaprender_el_cbc}{Fundamentos}, \linkUnidad{06_formas_lineales}{Formas lineales y dualidad} & \\
5 & Subespacios afines y lineales & \ref{cap:subespacios}. \nameref{cap:subespacios} & \linkUnidad{07a_cuna_determinante}{Cuña y determinante}, \linkUnidad{07b_nucleo_autovectores}{Núcleo y autovectores} & \\
6 & Subespacios invariantes & \ref{cap:autovectores}. \nameref{cap:autovectores} & \linkUnidad{07b_nucleo_autovectores}{Núcleo y autovectores}, \linkUnidad{07c_autovalores_3d_puntos_fijos}{Autovalores en 3D} & \\
7 & TL II: subespacios y la SVD & \ref{cap:svd}. \nameref{cap:svd} & \linkUnidad{09_cambio_de_base_svd}{Cambio de base, SVD e inversas} & \\
8 & Cambios de base & \ref{cap:cambio-base}. \nameref{cap:cambio-base} & \linkUnidad{08_cambio_de_base}{Cambio de base y diagonalización} & \\
9 & Repaso de análisis & \ref{cap:derivada}. \nameref{cap:derivada} & \linkUnidad{00_desaprender_el_cbc}{Fundamentos}, \linkUnidad{10_jacobiano_primera_forma}{Jacobiano y primera forma} & \\
10 & Inversas & \ref{cap:inversas}. \nameref{cap:inversas} & \linkUnidad{09_cambio_de_base_svd}{Cambio de base, SVD e inversas} & \\
11 & Cuatro problemas & \ref{cap:cuatro-problemas}. \nameref{cap:cuatro-problemas} & \linkUnidad{11_cuatro_problemas}{Cuatro problemas} & En preparación\\
12 & Curvatura & \ref{cap:curvatura}. \nameref{cap:curvatura} & \linkUnidad{12_curvatura}{Curvatura} & En preparación\\
13 & Geometría del elipsoide & \ref{cap:elipsoide}. \nameref{cap:elipsoide} & \linkUnidad{13_elipsoide}{Geometría del elipsoide}, \linkUnidad{15_transformaciones_conformes}{Transformaciones conformes} & En preparación\\
14 & Trigonometría esférica & \ref{cap:esfera}. \nameref{cap:esfera} & \linkUnidad{14_trigonometria_esferica}{Trigonometría esférica}, \linkUnidad{15_transformaciones_conformes}{Transformaciones conformes} & En preparación\\
--- & Transformaciones proyectivas & \ref{cap:proyectivas}. \nameref{cap:proyectivas} & \linkUnidad{05_proyectivas}{Proyectivas} & Tema de clase extra, dictado en dos fechas en 2026\\
--- & El terreno como matriz & \ref{cap:dem}. \nameref{cap:dem} & \linkUnidad{09_cambio_de_base_svd}{Cambio de base, SVD e inversas} & \\
--- & Transformaciones conformes & \ref{cap:conformes}. \nameref{cap:conformes} & \linkUnidad{15_transformaciones_conformes}{Transformaciones conformes} & \\
--- & Marcos de referencia, grillas y remuestreo & \ref{cap:marcos}. \nameref{cap:marcos} & --- & \\
\end{longtable}
}
